%-------------------- Abstract ---------------------------------------
\chapter*{Abstract}
\addcontentsline{toc}{chapter}{Abstract}
\thispagestyle{frontmatter}
Under India's E-Waste (Management) Rules, 2022, a producer of electrical and electronic equipment
must meet a yearly Extended Producer Responsibility (EPR) obligation by acquiring certificates
generated by registered recyclers, and a compliance officer must show, record by record, which
part of the obligation is actually supported by evidence. The records are heterogeneous, the
rules are temporal, and an error in either direction has regulatory consequences, so the work is
an investigation rather than a single lookup. This report presents ReCycleGuard, a local,
human-reviewed assistant for that investigation. Original files are stored immutably, extracted
with page-level locators, normalised with exact decimal arithmetic and reconciled by deterministic
tools. A bounded LangGraph agent plans an action, executes an audited tool, assesses the result
and either continues or stops; a deterministic gate accepts a model's finding only when computed
evidence supports it, and a person reviews the resulting working paper. The system was evaluated
on 48 synthetic dossiers covering sixteen fault families, with hidden answers kept outside the
runtime, using a locally served Qwen3.5-9B model. With the deterministic checks in place, the
single-agent loop succeeded on 32 of 32 development trials, but model-free controls scored the
same, and with the checks switched off no model variant recovered a judgement fault. A
pre-registered multi-agent design with a retrieval agent and a classification agent was then
tested across three prompt versions: no multi-agent arm passed the no-regression gate, and every
arm using the retrieval agent had 25 to 38 accepted false findings in 32 trials. On 24 sealed held-out
dossiers, run once as declared, the fixed and adaptive loops and both controls succeeded on every
trial (48 of 48 and 24 of 24) and the model contributed no accepted finding; the adaptive loop used
2.4 model calls per trial against 6.1. The main
limitation is that the dossiers are synthetic and were built by the same team that built and
evaluated the system.

\medskip
\noindent\textbf{Keywords:} agentic AI, extended producer responsibility, e-waste, tool use,
retrieval-augmented generation, human-in-the-loop, pre-registered evaluation.

%-------------------- Lists ------------------------------------------
\tableofcontents
\listoffigures
\addcontentsline{toc}{chapter}{List of Figures}
\listoftables
\addcontentsline{toc}{chapter}{List of Tables}

\chapter*{List of Abbreviations}
\addcontentsline{toc}{chapter}{List of Abbreviations}
\thispagestyle{frontmatter}
\begin{tabular}{@{}p{2.8cm}l@{}}
API & Application Programming Interface \\
BM25 & Best Matching 25 (lexical ranking function) \\
CAS & Compare-and-Swap \\
CI & Confidence Interval \\
CPCB & Central Pollution Control Board \\
EEE & Electrical and Electronic Equipment \\
EPR & Extended Producer Responsibility \\
FN & False Negative \\
FP & False Positive \\
FTS5 & SQLite Full-Text Search, version 5 \\
GPU & Graphics Processing Unit \\
HITL & Human-in-the-Loop \\
LLM & Large Language Model \\
MG1, MG2 & Model Gates 1 and 2 (the project's model experiments) \\
MoEFCC & Ministry of Environment, Forest and Climate Change \\
RAG & Retrieval-Augmented Generation \\
ReAct & Reasoning and Acting \\
SC & Scripted Control (model-free baseline) \\
TP & True Positive \\
VRAM & Video Random-Access Memory \\
\end{tabular}
